\part{Aktueller Stand}

\chapter{Was der Zufall mitbringt}

\section{Die Lücke, geschlossen}

\begin{merksatz}
\boxtitel{Was die anderen Bände nun haben}
\begin{center}
{\small
\renewcommand{\arraystretch}{1.5}
\begin{tabular}{@{}P{1.8cm}P{3.0cm}P{3.4cm}@{}}
\toprule
\textbf{Band} & \textbf{gebrauchte Größe} & \textbf{Kapitel hier}\\
\midrule
$16$ & Irrfahrten auf Graphen & $1$--$3$\\
$17$ & sequenzielles Testen & $9$\\
$18$ & Ausbruch oder Verlöschen & $6$\\
$22$ & Konvergenz des Gradientenabstiegs & $9$\\
$23$ & Filtration, Stoppzeiten & $7$--$8$\\
$24$ & Markow-Entscheidungen & $2$--$3$\\
\bottomrule
\end{tabular}}
\end{center}

\begin{center}
\begin{minipage}{0.9\textwidth}
\textbf{Sechs Bände hatten den Begriff gebraucht. Er steht nun auf eigenen Füßen.}
\end{minipage}
\end{center}

\begin{center}
\begin{minipage}{0.88\textwidth}
\kursiv{\textbf{Die dritte Zeile ist die folgenreichste.} Band 18 beschrieb Epidemien durch Differentialgleichungen und konnte deshalb nur sagen, was im Mittel geschieht. Kapitel~6 zeigt, dass bei $R_0 = 2$ jeder fünfte eingeschleppte Fall folgenlos bleibt --- und dass bei starker Überstreuung vier von fünf folgenlos bleiben. Beides ist mit denselben Mittelwerten vereinbar und führt zu verschiedenen Maßnahmen.}
\end{minipage}
\end{center}
\end{merksatz}

\section{Vier Muster}

\begin{merksatz}
\boxtitel{Muster 1 --- die Anschauung trügt}
\begin{center}
{\small
\renewcommand{\arraystretch}{1.48}
\begin{tabular}{@{}P{1.4cm}P{7.0cm}@{}}
\toprule
$1$ & Beim fairen Spiel führt meist einer fast durchgehend --- Ausgeglichenheit ist der seltenste Verlauf.\\
$4$ & Wer zufällig eintrifft, erlebt doppelt so lange Lücken wie der Durchschnitt.\\
$8$ & Auf KKK wartet man vierzehn Würfe, auf KZK zehn --- bei gleicher Häufigkeit.\\
$12$ & Nach dreißig Jahren liegen $71$ Prozent aller Verläufe unter dem Erwartungswert.\\
\bottomrule
\end{tabular}}
\end{center}

\begin{center}
\begin{minipage}{0.88\textwidth}
\kursiv{\textbf{Jedes dieser Ergebnisse widerspricht einer plausiblen Erwartung}, und keines beruht auf einem Trick. Sie alle folgen daraus, dass man beim Zufall über die Zeit nicht mit Mittelwerten allein auskommt --- die Verteilung der Pfade ist eine andere Sache als die Verteilung der Endstände.}
\end{minipage}
\end{center}
\end{merksatz}

\begin{merksatz}
\boxtitel{Muster 2 --- der Mittelwert beschreibt niemanden}
\begin{center}
{\small
\renewcommand{\arraystretch}{1.48}
\begin{tabular}{@{}P{1.4cm}P{7.0cm}@{}}
\toprule
$3$ & Sechzig Prozent stabil: jeder Patient zu sechzig Prozent --- oder sechzig Prozent der Patienten immer.\\
$7$ & Die bedingte Erwartung ist ein Mittelwert, kein typischer Wert.\\
$9$ & Die Urne konvergiert gegen einen zufälligen Grenzwert, nie gegen den Erwartungswert.\\
$12$ & Erwartungswert und Median der Anlage: $11{,}0$ gegen $6{,}1$.\\
\bottomrule
\end{tabular}}
\end{center}

\begin{center}
\begin{minipage}{0.88\textwidth}
\kursiv{\textbf{Bei multiplikativen Vorgängen ist der Mittelwert fast immer die falsche Kennzahl.} Er wird von wenigen großen Verläufen getragen, während die Mehrheit darunter bleibt --- bei Vermögen, Klongrößen, Konzentrationen und Populationen gleichermaßen. Der Median ist hier nicht die robustere, sondern die zutreffendere Größe.}
\end{minipage}
\end{center}
\end{merksatz}

\begin{merksatz}
\boxtitel{Muster 3 --- Schwankung kostet}
\begin{center}
{\small
\renewcommand{\arraystretch}{1.48}
\begin{tabular}{@{}P{1.4cm}P{7.0cm}@{}}
\toprule
$4$ & Kleine Einheiten brauchen anteilig mehr Reserve --- die Wurzelregel.\\
$5$ & Die letzten fünf Prozent Auslastung kosten mehr Wartezeit als die ersten neunzig.\\
$12$ & Der typische Verlauf wächst mit $\mu - \sigma^{2}/2$, nicht mit $\mu$.\\
$13$ & Die Nachbildung einer Option ist nur im Mittel genau.\\
\bottomrule
\end{tabular}}
\end{center}

\begin{center}
\begin{minipage}{0.88\textwidth}
\kursiv{\textbf{Schwankung ist nicht nur ein Risiko, sondern ein Abzug} --- sie verschlechtert das Ergebnis auch dann, wenn niemand Risiko scheut. Das gilt für Wartezeiten, für Wachstum und für Kapazitäten, und es ist der mathematische Kern jedes Arguments für Puffer, Bündelung und Vergleichmäßigung.}
\end{minipage}
\end{center}
\end{merksatz}

\begin{merksatz}
\boxtitel{Muster 4 --- was man wann weiß}
\begin{center}
{\small
\renewcommand{\arraystretch}{1.48}
\begin{tabular}{@{}P{1.4cm}P{7.0cm}@{}}
\toprule
$7$ & Die Filtration ist die Fassung von \enquote{was zu diesem Zeitpunkt bekannt ist}.\\
$8$ & Eine Stoppzeit darf nicht in die Zukunft sehen.\\
$11$ & Nur die linke Auswertung macht den Integranden zur Handlungsanweisung.\\
$14$ & Der Filter schätzt den Zustand aus Vergangenem, nicht aus Kommendem.\\
\bottomrule
\end{tabular}}
\end{center}

\begin{center}
\begin{minipage}{0.88\textwidth}
\kursiv{\textbf{Dasselbe Prinzip trug in Band 22 die Warnung vor Datenlecks und in Band 23 die unsterbliche Zeit.} Hier ist es kein methodischer Ratschlag, sondern ein Begriff --- und der Unterschied zwischen dem Itō- und dem Stratonowitsch-Integral zeigt, dass er Folgen hat, die man ausrechnen kann: genau $T/2$.}
\end{minipage}
\end{center}
\end{merksatz}

\section{Was offen ist}

\begin{merksatz}
\boxtitel{Vier Fragen}
\begin{center}
{\small
\renewcommand{\arraystretch}{1.55}
\begin{tabular}{@{}P{3.0cm}P{5.6cm}@{}}
\toprule
\textbf{Selbstvermeidung} & Wie weit entfernt sich eine Irrfahrt ohne Wiederkehr? In drei Dimensionen unbekannt.\\
\textbf{Rauhe Pfade} & Eine pfadweise Theorie stochastischer Gleichungen ohne Wahrscheinlichkeit.\\
\textbf{Mischzeiten} & Wann hat eine Simulation eingeschwungen? Im Einzelfall meist unbeantwortet.\\
\textbf{Modellwahl} & Aus einem Verlauf auf den erzeugenden Prozess schließen.\\
\bottomrule
\end{tabular}}
\end{center}

\begin{center}
\begin{minipage}{0.88\textwidth}
\kursiv{Die letzte Frage verbindet diesen Band mit Band 23. \textbf{Ein beobachteter Verlauf ist mit sehr verschiedenen Mechanismen vereinbar} --- eine Häufung kann aus Ansteckung, aus schwankender Rate oder aus Überstreuung stammen, und alle drei erzeugen ähnliche Bilder. Die Prozesstheorie liefert die Modelle; welches zutrifft, entscheidet sie nicht.}
\end{minipage}
\end{center}
\end{merksatz}

\section{Sechs Sätze, die bleiben}

\begin{merksatz}
\boxtitel{Die Bilanz}
\begin{center}
{\small
\renewcommand{\arraystretch}{1.55}
\begin{tabular}{@{}P{0.7cm}P{7.7cm}@{}}
\toprule
$1$ & Eine Irrfahrt kehrt in einer und zwei Dimensionen sicher zurück, in drei nur mit Wahrscheinlichkeit $0{,}34$.\\
$2$ & Eine irreduzible, aperiodische Kette vergisst ihren Ausgangspunkt --- mit der Rate des zweitgrößten Eigenwerts.\\
$3$ & Drei harmlose Annahmen erzwingen den Poisson-Prozess samt exponentiellen Wartezeiten.\\
$4$ & Kein Anhaltekriterium macht aus einem fairen Spiel ein günstiges.\\
$5$ & Ein Martingal, das im Mittel beschränkt bleibt, konvergiert fast sicher --- gegen einen zufälligen Wert.\\
$6$ & Für zerknitterte Pfade gilt $(\dd W)^{2} = \dd t$; jede Kettenregel trägt deshalb einen Zusatzterm.\\
\bottomrule
\end{tabular}}
\end{center}
\end{merksatz}

\begin{merksatz}
\boxtitel{Ein Schlusswort}
\begin{center}
\begin{minipage}{0.9\textwidth}
Dieser Band begann mit einer Münze, die zweitausendmal geworfen wird, und endete bei Nervenzellen, Genen und Optionen. Dazwischen lag ein einziger Gedanke: Zufall, der weitergeht, ist etwas anderes als Zufall, der einmal eintritt.

Der Unterschied liegt nicht in den Verteilungen --- die sind meist dieselben wie in Band 4 --- sondern in den Pfaden. Ein Pfad hat eine Gestalt: Er steigt, kehrt zurück, bleibt lange auf einer Seite, erreicht Höchststände, stirbt aus oder erstarrt. Über diese Gestalt sagt die Verteilung des Endstands nichts, und fast alle Fragen, die man in der Anwendung stellt, betreffen die Gestalt.

Daraus folgt eine Haltung, die über die Mathematik hinausreicht. Wer einen Verlauf sieht --- einen Kurs, eine Fieberkurve, eine Reihe von Zwischenfällen --- sieht zwangsläufig Muster. Die Sätze dieses Bandes sagen, welche davon auch reiner Zufall erzeugt hätte: lange einseitige Strecken, Häufungen, scheinbare Trends, plötzliche Sprünge. Sie beweisen nicht, dass keine Ursache vorliegt. Sie sagen nur, dass man ohne Rechnung nicht entscheiden kann, ob eine nötig ist.
\end{minipage}
\end{center}

\begin{center}
\begin{minipage}{0.88\textwidth}
\kursiv{Damit steht dieser Band zwischen den Grundlagen und den Anwendungen. Die Wahrscheinlichkeitsrechnung sagt, was einmal geschehen kann. Die Statistik sagt, was die Daten hergeben. Das Lernen sagt, was auf den nächsten Fall zutrifft, die Kausalität, was sich ändert, wenn man handelt, und die Entscheidungstheorie, wie man wählt. \textbf{Und die Prozesstheorie sagt, wie es weitergeht --- und was daran schon deshalb so aussieht, weil die Zeit vergeht.}}
\end{minipage}
\end{center}
\end{merksatz}
